% SOURCE_AUTHORITY: sga5_fr_workpass.tex, lines 12474--12495 (LF-normalized unit).
% SOURCE_SHA256: 791F4EFFC5E02832D5D77ED03518C8156D6F07E4C8238B03545DB93D883FBB28
\begin{statement}{Proposição 3.8}
Sejam $G$ um grupo finito, $P^\bullet$ um complexo limitado de $\Lambda[G]$-módulos cujas componentes são todas projetivas de tipo finito, e $M^\bullet$ um complexo de $A[G]$-módulos. Nestas condições, existem isomorfismos canônicos de complexos de $A$-módulos:
\[
\Hom_{\Lambda[G]}^\bullet(P^\bullet,M^\bullet)
\simeq \Hom_\Lambda^\bullet(P^\bullet,M^\bullet)^G
\simeq (\check P^\bullet\otimes_\Lambda M^\bullet)^G
\simeq \check P^\bullet\otimes_{\Lambda[G]}M^\bullet.
\tag{3.8.1}
\]
\end{statement}

A demonstração utiliza os isomorfismos bem conhecidos para o caso dos módulos, e é deixada ao leitor. Observe-se que, no segundo e no terceiro termo, as operações de $G$ sobre $\Hom_\Lambda^\bullet(P^\bullet,M^\bullet)$, resp. $\check P^\bullet\otimes_\Lambda M^\bullet$, se definem por transporte de estrutura a partir de suas operações sobre $P^\bullet$ e $M^\bullet$, resp. $\check P^\bullet$ e $M^\bullet$.

\begin{statement}{Corolário 3.9}
Se $P^\bullet\in \operatorname{Ob}(D(\Lambda[G])_{\parf})$ e $M^\bullet\in D^+(A[G])$, então se têm na categoria derivada $D(A)$ os isomorfismos seguintes:
\[
\RHom_{\Lambda[G]}(P^\bullet,M^\bullet)
=\R\Gamma^G(\check P^\bullet\otimes_\Lambda^{\mathbf L}M^\bullet)
\simeq \check P^\bullet\otimes_{\Lambda[G]}M^\bullet.
\tag{3.9.1}
\]
\end{statement}
